%------------------------------------------------------------------------------%
\documentclass[fleqn,utf8,aspectratio=169,12pt]{beamer}

\usepackage{calculo_varias_variaveis-1}

\title{Máximos e Mínimos}

%------------------------------------------------------------------------------%
\begin{document}

\addtitle

%------------------------------------------------------------------------------%
\section{Extremos locais}

%------------------------------------------------------------------------------%
\begin{frame}{Extremos locais}
  Seja $f(x, y)$ definida em uma região $R$ que contém o ponto $(a, b)$

  \pause
  \begin{itemize}[<+->]
    \item \(f(a, b)\) é um valor \emph{máximo local} de $f$ se
          \[
            \emph{f(a, b) \;\geq\; f(x, y)}
          \]
          para \((x, y)\) do domínio em um disco aberto centrado em \((a, b)\)
          \\[5mm]
    \item \(f(a, b)\) é um valor \emph{mínimo local} de $f$ se
          \[
            \emph{f(a, b) \;\leq\; f(x, y)}
          \]
          para \((x, y)\) do domínio em um disco aberto centrado em \((a, b)\)
  \end{itemize}
\end{frame}

%------------------------------------------------------------------------------%
\begin{frame}{Teste da derivada de primeira ordem}
  Se $f(x, y)$ tiver um máximo ou mínimo local
  em um ponto interior \((a, b)\) \\
  do seu domínio

  \pause
  \medskip
  e se as derivadas de primeira ordem existirem em \((a, b)\), então

  \pause
  \medskip
  \[
    \emph{f_x(a, b) = 0}
    \qquad
    \pause
    \emph{f_y(a, b) = 0}
  \]
\end{frame}

%------------------------------------------------------------------------------%
\begin{frame}{Ponto crítico}

  Um \emph{Ponto Crítico} de uma função \(f(x, y)\)

  \pause
  \smallskip
  é um ponto $(a, b)$, no \emph{interior do domínio} de $f$, onde

  \pause
  \bigskip
  \emph{\(f_x(a, b) = 0\)} \; e \; \emph{\(f_y(a, b) =0\)}
  \pause
  \tabto{70mm}
  \(\nabla f(a, b) = 0\)

  \pause
  ou \par
  \(f_x(a, b)\) \; ou \; \(f_y(a, b)\) \; \emph{não exista}
  \pause
  \tabto{70mm}
  \(\nabla f(a, b)\) não existe

\end{frame}

%------------------------------------------------------------------------------%
\begin{frame}{Ponto de sela}
  Uma função diferenciável \(f(x, y)\) possui um \emph{ponto de sela} em \((a, b)\)

  \pause
  se, em todo disco aberto centrado em \((a, b)\),

  \pause
  existirem pontos \((x_1, y_1)\) e \((x_2, y_2)\) onde

  \pause
  \medskip
  \[
    \emph{f(x_1, y_1) \;>\; f(a, b)}
  \]
  \[
    \pause
    \emph{f(x_2, y_2) \;<\; f(a, b)}
  \]
\end{frame}

%------------------------------------------------------------------------------%
\include{J-01-Maximos_minimos-ex1}

%------------------------------------------------------------------------------%
\section{Teste da derivada segunda}

%------------------------------------------------------------------------------%
\begin{frame}{Hessiana}
  \emph{Matriz Hessiana}
  \[
    \pause
    \arraycolsep=2mm
    \def\arraystretch{1.2}
    H(x, y)
    \pause
      \;=\; \begin{bmatrix}
      \null\;f_{xx}(x, y)\;\null & \null\;f_{xy}(x, y)\;\null \\
      f_{yx}(x, y) & f_{yy}(x, y)
    \end{bmatrix}
  \]

  \pause
  \bigskip
  \emph{Discriminante (ou Hessiano)}
  \[
    \pause
    D(x, y)
    \pause
    \;=\; \det H(x, y)
    \pause
    \;=\; f_{xx}(x, y)\,f_{yy}(x, y) \;-\; f_{xy}^2(x, y)
  \]
\end{frame}

%------------------------------------------------------------------------------%
\begin{frame}{Teste da derivada segunda}
  Seja $f(x, y)$ diferenciável
  \pause
  e $(a, b)$ um ponto crítico onde
  \emph{\(
    f_x(a, b) = f_y(a, b) = 0
  \)}

  \pause
  \begin{enumerate}[<+->]
    \setlength{\itemsep}{4mm}
    \item se \emph{\(\det H(a, b) = 0\)} \tabto{40mm}
      o teste é inconclusivo
    \item se \emph{\(\det H(a, b) < 0\)} \tabto{40mm}
      $f$ tem um \emph{ponto de sela} em $(a, b)$
    \item se \emph{\(\det H(a, b) > 0\)} \tabto{40mm}
      testamos a derivada segunda em $x$ (ou $y$) \\[3mm]
      \begin{enumerate}
        \setlength{\itemsep}{3mm}
        \item \normalsize\; se \emph{\(f_{xx}(a, b) < 0\)} \tabto{35mm}
          $f$ tem um \emph{máximo local} em $(a, b)$
        \item \normalsize\; se \emph{\(f_{xx}(a, b) > 0\)} \tabto{35mm}
          $f$ tem um \emph{mínimo local} em $(a, b)$
      \end{enumerate}
  \end{enumerate}
\end{frame}

%------------------------------------------------------------------------------%
\section{Exemplos}

\include{ex/J-01-Maximos_minimos-ex2}
\include{ex/J-01-Maximos_minimos-ex3}
\include{ex/J-01-Maximos_minimos-ex4}
\include{ex/J-01-Maximos_minimos-ex5}
\include{ex/J-01-Maximos_minimos-ex6}
%------------------------------------------------------------------------------%
\begin{question}
  Considerando a função
  \(
    f(x, y) = x^3 + 3xy + y^3
  \)

  \smallskip
  a) Calcule o gradiente de $f$

  \smallskip
  b) Calcule a hessiana de $f$

  \smallskip
  c) Encontre todos os pontos críticos de $f$

  \smallskip
  d) Classifique cada ponto crítico de $f$
\end{question}

%------------------------------------------------------------------------------%
\begin{resolution}{a)\; Gradiente de $f$}
  Derivadas parciais de $f$

  \[
    \pause
    \D{f}{x}
    \pause
    \;=\; \D{}{x}\left[x^3 + 3xy + y^3\right]
    \pause
    \;=\; 3x^2 + 3y
  \]

  \[
    \pause
    \D{f}{y}
    \pause
    \;=\; \D{}{y}\left[x^3 + 3xy + y^3\right]
    \pause
    \;=\; 3x + 3y^2
  \]

  \pause
  então
  \[
    \pause
    \nabla f = \vetor{3x^2 + 3y}{3x + 3y^2}
  \]
\end{resolution}

%------------------------------------------------------------------------------%
\begin{resolution}{b)\; Hessiana de $f$}
  Derivadas parciais de segunda ordem de $f$
  \[
    \pause
    \DS{f}{x}
    \pause
    \;=\; \D{}{x}\left[3x^2 + 3y\right]
    \pause
    \;=\; 6x
  \]
  \[
    \pause
    \D{f}{y}
    \pause
    \;=\; \D{}{y}\left[3x + 3y^2\right]
    \pause
    \;=\; 6y
  \]
  \[
    \pause
    \DM{f}{x}{y}
    \pause
    \;=\; \D{}{x}\D{f}{y}
    \pause
    \;=\; \D{}{x}\left[3x + 3y^2\right]
    \pause
    \;=\; 3
  \]

  \pause
  então
  \[
    \pause
    H = \left(\begin{array}{cc}
      6x & 3 \\
      3  & 6y
    \end{array}\right)
  \]
\end{resolution}

%------------------------------------------------------------------------------%
\begin{resolution}{c)\; Pontos Críticos de $f$}
  A função é um polinômio, então possui derivadas em todos os pontos do plano

  \pause
  Os pontos críticos são apenas onde as derivadas parciais são zero,
  \(\nabla f = 0\)

  \[
    \pause
    3x^2 + 3y = 0
    \qquad\text{e}\qquad
    \pause
    3x + 3y^2 = 0
  \]
  \pause
  simplificando
  \[
    \pause
    x^2 + y = 0
    \qquad\text{e}\qquad
    \pause
    x + y^2 = 0
  \]
\end{resolution}

%------------------------------------------------------------------------------%
\begin{resolution}{c)\; Pontos Críticos de $f$}
  \onslide<+->{Isolando $y$ na primeira equação}
  \begin{align*}
    \onslide<+->{x^2 + y & = 0} \\
    \onslide<+->{y & = -x^2}
  \end{align*}
  \onslide<+->{e substituindo na segunda}
  \begin{align*}
    \onslide<+->{x + y^2                 & = 0 \\}
    \onslide<+->{x + \left(-x^2\right)^2 & = 0 \\}
    \onslide<+->{x + x^4                 & = 0 \\}
    \onslide<+->{x(1+ x^3)               & = 0}
  \end{align*}
  \onslide<+->{As soluções dessa equação são $x=0$ ou $x-1$}
\end{resolution}

%------------------------------------------------------------------------------%
\begin{resolution}{c)\; Pontos Críticos de $f$}
  Como \(y = -x^2\)

  \pause
  \medskip
  se $x=0$, então $y=0$

  \pause
  \medskip
  se $x=-1$, então $y=-1$

  \pause
  \bigskip
  os pontos críticos são
  \[
    (x_1, y_1) = (0, 0)
    \pause
    \qquad\text{e}\qquad
    (x_2, y_2) = (-1, -1)
  \]
\end{resolution}

%------------------------------------------------------------------------------%
\begin{resolution}{d)\; Classificando os Pontos}
  Considerando o ponto \((x_1, y_1) = (0, 0)\)
  \[
    \pause
    f_{xx}(0, 0)
    \pause
    = \left(6x\right)\evalat{(0, 0)}{}
    \pause
    = 0
  \]
  \[
    \pause
    f_{yy}(0, 0)
    \pause
    = \left(6y\right)\evalat{(0, 0)}{}
    \pause
    = 0
  \]
  \[
    \pause
    f_{xy}(0, 0) = 3
  \]
  \[
    \pause
    D_1
    \pause
    \;=\; f_{xx}(0, 0)f_{yy}(0, 0) - f^2_{xy}(0, 0)
    \pause
    \;=\; 0 \times 0 - 3^2
    \pause
    \;=\; -9
    \pause
    \;<\; 0
  \]

  \pause
  Portanto, o ponto $(0, 0)$ é um ponto de sela.
\end{resolution}

%------------------------------------------------------------------------------%
\begin{resolution}{d)\; Classificando os Pontos}
  Considerando o ponto \((x_2, y_2) = (-1, -1)\)
  \[
    \pause
    f_{xx}(-1, -1)
    \pause
    = \left(6x\right)\evalat{(-1, -1)}{}
    \pause
    = -6
  \]
  \[
    \pause
    f_{yy}(-1, -1)
    \pause
    = \left(6y\right)\evalat{(-1, -1)}{}
    \pause
    = -6
  \]
  \[
    \pause
    f_{xy}(-1, -1) = 3
  \]
  \[
    \pause
    D_2
    \pause
    \;=\; f_{xx}(-1, -1)f_{yy}(-1, -1) - f^2_{xy}(-1, -1)
    \pause
    \;=\; (-6)(-6) - 3^2
    \pause
    \;=\; 36 - 9
    \pause
    \;=\; 25 > 0
  \]

  \pause
  Portanto, o ponto $(-1, -1)$ é um máximo ou mínimo local

  \pause
  Como \(f_{xx}(-1, -1) = -6 < 0\)

  \pause
  o ponto é um ponto de máximo local.
\end{resolution}

%------------------------------------------------------------------------------%

%------------------------------------------------------------------------------%
\begin{question}
  Considerando a função
  \(
    f(x, y) = 4xy -x^4 -y^4
  \)

  a) Calcule o gradiente de $f$

  b) Calcule a hessiana de $f$

  c) Encontre todos os pontos críticos de $f$

  d) Classifique cada ponto crítico de $f$
\end{question}

%------------------------------------------------------------------------------%
\begin{resolution}{a -- Gradiente de $f$}
  \[
    \D{f}{x}
    \pause
    \;=\; \D{}{x}\left[4xy -x^4 -y^4\right]
    \pause
    \;=\; 4y - 4x^3
  \]

  \pause
  \[
    \D{f}{y}
    \pause
    \;=\;  \D{}{y}\left[4xy -x^4 -y^4\right]
    \pause
    \;=\;  4x -4y^3
  \]

  \pause
  \[
    \text{então}\qquad
    \nabla f(x, y) \;=\; \vetor{4y - 4x^3}{4x -4y^3}
  \]
\end{resolution}

%------------------------------------------------------------------------------%
\begin{resolution}{b -- Hessiana de $f$}
  \[
    \DS{f}{x}
    \pause
    \;=\; \D{}{x}\left[4y - 4x^3\right]
    \pause
    \;=\; -12x^2
  \]

  \pause
  \[
    \DS{f}{y}
    \pause
    \;=\; \D{}{y}\left[4x -4y^3\right]
    \pause
    \;=\; -12y^2
  \]

  \pause
  \[
    \DM{f}{x}{y}
    \pause
    \;=\; \D{}{x}\D{f}{y}
    \pause
    \;=\; \D{}{x}\left[4x -4y^3\right]
    \pause
    \;=\; 4
  \]
  \pause
  \[
    \text{então}\qquad
    H(x, y)
    \;=\; \left(\begin{array}{cc}
      -12x^2 & 4      \\
      4      & -12y^2
    \end{array}\right)
  \]
\end{resolution}

%------------------------------------------------------------------------------%
\begin{resolution}{c -- Pontos críticos de $f$}
  A função é um polinômio, então possui derivadas em todos os pontos do plano

  \pause
  Assim os pontos críticos são apenas os pontos onde \(\nabla f = 0\)

  \pause
  \[
    4y - 4x^3 = 0
    \qquad\text{e}\qquad
    4x -4y^3 = 0
  \]

  \pause
  ou
  \[
    y - x^3 = 0
    \qquad\text{e}\qquad
    x - y^3 = 0
  \]
\end{resolution}

%------------------------------------------------------------------------------%
\begin{resolution}{c -- Resolvendo o sistema}
\begin{columns}[t]
\column{0.45\linewidth}
  \onslide<+->{Isolando $y$ na primeira equação}
  \begin{align*}
    \onslide<+->{y - x^3 & = 0   \\}
    \onslide<+->{y       & = x^3}
  \end{align*}
\column{0.55\linewidth}
  \onslide<+->{e substituindo na segunda}
  \begin{align*}
    \onslide<+->{x - y^3                & = 0 \\}
    \onslide<+->{x - \left(x^3\right)^3 & = 0 \\}
    \onslide<+->{x -x^9                 & = 0 \\}
    \onslide<+->{x\left(1 - x^8\right)  & = 0}
  \end{align*}
  \onslide<+->{As soluções dessa equação são\\$x=0$, $x=1$ ou $x-1$}
\end{columns}
\end{resolution}

%------------------------------------------------------------------------------%
\begin{resolution}{c -- Resolvendo o sistema}
  Como \(y=x^3\)

  \pause
  Se \(x=0\) temos \(y=0\)

  \pause
  Se \(x=1\) temos \(y=1\)

  \pause
  Se \(x=-1\) temos \(y=-1\)

  \pause
  \bigskip
  Os pontos críticos são
  \[
    (x_1, y_1) = (0, 0),
    \qquad
    (x_2, y_2) = (1, 1)
    \quad\text{e}\quad
    (x_3, y_3) = (-1, -1)
  \]
\end{resolution}

%------------------------------------------------------------------------------%
\begin{resolution}{d -- Classificando o ponto \((x_1, y_1) = (0, 0)\)}
\begin{columns}[t]
\column{0.45\linewidth}
  \[
    f_{xx}(0, 0)
    \pause
    \;=\; \left(-12x^2\right) \evalat{(0, 0)}{}
    \pause
    =\; 0
  \]
  \[
    \pause
    f_{yy}(0, 0)
    \pause
    \;=\; \left(-12y^2\right) \evalat{(0, 0)}{}
    \pause
    =\; 0
  \]
  \[
    \pause
    f_{xy}(0, 0)
    \pause
    \;=\; \left(4\right) \evalat{(0, 0)}{}
    \pause
    =\; 4
  \]
\column{0.55\linewidth}
  \pause
  \begin{align*}
    \onslide<+->{D_1}
    \onslide<+->{& = \det H(0, 0) \\}
    \onslide<+->{& = f_{xx}(0, 0)f_{yy}(0, 0) - f^2_{xy}(0, 0) \\}
    \onslide<+->{& = 0 \times 0 - 4^2 \\}
    \onslide<+->{& = -16 }
    \onslide<+->{< 0}
  \end{align*}
\end{columns}

  \bigskip
  \onslide<+->{O ponto $(0, 0)$ é um \emph{ponto de sela}}
\end{resolution}

%------------------------------------------------------------------------------%
\begin{resolution}{d -- Classificando o ponto \((x_2, y_2) = (1, 1)\)}
\begin{columns}[t]
\column{0.45\linewidth}
  \[
    f_{xx}(1, 1)
    \pause
    \;=\; \left(-12x^2\right) \evalat{(1, 1)}{}
    \pause
    =\; -12
  \]
  \[
    \pause
    f_{yy}(1, 1)
    \pause
    \;=\; \left(-12y^2\right) \evalat{(1, 1)}{}
    \pause
    =\; -12
  \]
  \[
    \pause
    f_{xy}(1, 1)
    \pause
    \;=\; \left(-12x^2\right) \evalat{(1, 1)}{}
    \pause
    =\; 4
  \]
\column{0.55\linewidth}
  \pause
  \begin{align*}
    \onslide<+->{D_2}
    \onslide<+->{& = \det H(1, 1) \\}
    \onslide<+->{& = f_{xx}(1, 1)f_{yy}(1, 1) - f^2_{xy}(1, 1) \\}
    \onslide<+->{& = (-12)(-12) - 4^2 \\}
    \onslide<+->{& = 144 -16 \\}
    \onslide<+->{& = 128 }
    \onslide<+->{> 0}
  \end{align*}
\end{columns}

  \bigskip
  \onslide<+->{O ponto $(1, 1)$ é um extremo local}

  \onslide<+->{Como \(f_{xx}(1, 1) = -12 < 0\)}
  \onslide<+->{o ponto é um \emph{máximo local}}
\end{resolution}

%------------------------------------------------------------------------------%
\begin{resolution}{d -- Classificando o ponto \((x_3, y_3) = (-1, -1)\)}
\begin{columns}[t]
\column{0.5\linewidth}
  \[
    f_{xx}(-1, -1)
    \pause
    = \left(-12x^2\right) \evalat{(-1, -1)}{}
    \pause
    \!\!= -12
  \]
  \[
    \pause
    f_{yy}(-1, -1)
    \pause
    = \left(-12x^2\right) \evalat{(-1, -1)}{}
    \pause
    \!\!= -12
  \]
  \[
    \pause
    f_{xy}(-1, -1)
    \pause
    = \left(-12x^2\right) \evalat{(-1, -1)}{}
    \pause
    \!\!= 4
  \]
\column{0.5\linewidth}
  \pause
  \begin{align*}
    \!
    \onslide<+->{D_3}
    \onslide<+->{& = \det H(-1, -1) \\}
    \onslide<+->{& = f_{xx}(-1, -1)f_{yy}(-1, -1) - f^2_{xy}(-1, -1) \\}
    \onslide<+->{& = (-12)(-12) - 4^2 \\}
    \onslide<+->{& = 144 -16 \\}
    \onslide<+->{& = 128 }
    \onslide<+->{> 0}
  \end{align*}
\end{columns}

  \bigskip
  \onslide<+->{O ponto $(-1, -1)$ é um extremo local}

  \onslide<+->{Como \(f_{xx}(-1, -1) = -12 < 0\)}
  \onslide<+->{o ponto é um \emph{máximo local}}
\end{resolution}

%------------------------------------------------------------------------------%

%------------------------------------------------------------------------------%
\begin{question}
  Considerando a função
  \(
    f(x, y) = x^4 + y^4 + 4xy
  \)

  a) Calcule o gradiente de $f$

  b) Calcule a hessiana de $f$

  c) Encontre todos os pontos críticos de $f$

  d) Classifique cada ponto crítico de $f$
\end{question}

%------------------------------------------------------------------------------%
\begin{resolution}{a -- Gradiente de $f$}
  \[
    \D{f}{x}
    \pause
    \;=\; \D{}{x}\left[x^4 + y^4 + 4xy\right]
    \pause
    \;=\; 4x^3 + 4y
  \]

  \pause
  \[
    \D{f}{y}
    \pause
    \;=\; \D{}{y}\left[x^4 + y^4 + 4xy\right]
    \pause
    \;=\; 4y^3 + 4x
  \]

  \pause
  \[
    \text{então}\qquad
    \nabla f = \vetor{4x^3 + 4y}{4y^3 + 4x}
  \]
\end{resolution}

%------------------------------------------------------------------------------%
\begin{resolution}{b -- Hessiana de $f$}
  \[
    \DS{f}{x}
    \pause
    \;=\; \D{}{x}\left[4x^3 + 4y\right]
    \pause
    \;=\; 12x^2
  \]

  \pause
  \[
    \DS{f}{y}
    \pause
    \;=\; \D{}{y}\left[4y^3 + 4x\right]
    \pause
    \;=\; 12y^2
  \]

  \pause
  \[
    \DM{f}{x}{y}
    \pause
    \;=\; \D{}{x}\D{f}{y}
    \pause
    \;=\; \D{}{x}\left[4y^3 + 4x\right]
    \pause
    \;=\; 4
  \]

  \pause
  \[
    \text{então}\qquad
    H = \left(\begin{array}{cc}
      12x^2 & 4     \\
      4     & 12y^2
    \end{array}\right)
  \]
\end{resolution}

%------------------------------------------------------------------------------%
\begin{resolution}{c -- Pontos críticos}
  A função é um polinômio, então possui derivadas em todos os pontos do plano

  \pause
  Assim os pontos críticos são apenas os pontos onde \(\nabla f = 0\)

  \pause
  \[
    4x^3 + 4y = 0
    \qquad\text{e}\qquad
    4y^3 + 4x = 0
  \]
  \pause
  ou
  \[
    x^3 + y = 0
    \qquad\text{e}\qquad
    y^3 + x = 0
  \]
\end{resolution}

%------------------------------------------------------------------------------%
\begin{resolution}{c -- Resolvendo o sistema}
  Isolando $y$ na primeira equação, \(y = -x^3\), e substituindo na segunda, temos
  \begin{align*}
    y^3 + x                 & = 0 \\
    \left(-x^3\right)^3 + x & = 0 \\
    -x^9 + x                & = 0 \\
    x^9 - x                 & = 0 \\
    x(x^8 - 1)              & = 0
  \end{align*}
  As soluções dessa equação são $x=0$, $x=1$ ou $x-1$.
  Se $x=0$ temos $y=0$,
  se $x=1$ temos $y=-1$ e
  se $x=-1$ temos $y=1$.
  Portanto, os pontos críticos são
  \[
    (x_1, y_1) = (0, 0),
    \qquad\qquad
    (x_2, y_2) = (1, -1)
    \qquad\text{e}\qquad
    (x_3, y_3) = (-1, 1)
  \]
\end{resolution}

%------------------------------------------------------------------------------%
\begin{resolution}{d -- Classificando o ponto \((x_1, y_1) = (0, 0)\)}
  \[
  f_{xx}(0, 0) = 0
  \]
  \[
  f_{yy}(0, 0) = 0
  \]
  \[
  f_{xy}(0, 0) = 4
  \]
  \[
  D_1
  = f_{xx}(0, 0)f_{yy}(0, 0) - f^2_{xy}(0, 0)
  = 0 \times 0 - 4^2
  = -16 < 0
  \]
  Portanto, o ponto $(0, 0)$ é um ponto de sela.
\end{resolution}

%------------------------------------------------------------------------------%
\begin{resolution}{d -- Classificando o ponto \((x_2, y_2) = (1, -1)\)}
  \[
    f_{xx}(1, -1) = 12
  \]
  \[
    f_{yy}(1, -1) = 12
  \]
  \[
    f_{xy}(1, -1) = 4
  \]
  \[
    D_2
    = f_{xx}(1, -1)f_{yy}(1, -1) - f^2_{xy}(1, -1)
    = 12 \times 12 - 4^2
    = 144 -16
    = 128 > 0
  \]
  Portanto, o ponto $(1, -1)$ é um máximo ou mínimo local.
  Como \(f_{xx}(1, -1) = 12 > 0\)
  o ponto é um ponto de mínimo local.
\end{resolution}

%------------------------------------------------------------------------------%
\begin{resolution}{d -- Classificando o ponto \((x_3, y_3) = (-1, 1)\)}
  \[
    f_{xx}(-1, 1) = 12
  \]
  \[
    f_{yy}(-1, 1) = 12
  \]
  \[
    f_{xy}(-1, 1) = 4
  \]
  \[
    D_3
    = f_{xx}(-1, 1)f_{yy}(-1, 1) - f^2_{xy}(-1, 1)
    = 12 \times 12 - 4^2
    = 144 -16
    = 128 > 0
  \]
  Portanto, o ponto $(-1, 1)$ é um máximo ou mínimo local.
  Como \(f_{xx}(-1, 1) = 12 > 0\)
  o ponto é um ponto de mínimo local.
\end{resolution}

%------------------------------------------------------------------------------%


%------------------------------------------------------------------------------%
\section{Lista mínima}

%------------------------------------------------------------------------------%
\begin{frame}{Lista mínima}
  Cálculo Vol. 2 do Thomas 12\textsuperscript{a} ed. -- Seção 14.7
  \begin{enumerate}
    \item Estudar o texto da seção
    \item Resolver os exercícios: 2, 9, 11, 21, 23, 24, 25, 27
  \end{enumerate}

  \vfill
  Atenção: A prova é baseada no livro, não nas apresentações
\end{frame}

%------------------------------------------------------------------------------%
\end{document}
%------------------------------------------------------------------------------%
